% =========================================================
%  references.tex
%  Rujukan yang dipakai modul.
% =========================================================

\clearpage
\phantomsection
\addcontentsline{toc}{section}{Rujukan}
\section*{Rujukan}

\begin{thebibliography}{9}
\setlength{\itemsep}{0.35em}

\bibitem{bukuajar}
Cecep Suwanda. \emph{Pengantar Pemrograman Game: Buku Ajar Berbasis
\istilah{Outcome-Based Education} (OBE)}. Universitas Bale Bandung, 2026.
Bab III--XI menjadi sumber konsep modul ini.

\bibitem{rps}
Cecep Suwanda. \emph{Rencana Pembelajaran Semester \NamaMK\ (\KodeMK)}.
Program Studi \NamaProdi, Fakultas \NamaFakultas, Universitas Bale Bandung,
Semester \SemesterMK\ \TahunAkademik.

\bibitem{pygame}
Pygame Community. \emph{Pygame Documentation}. \url{https://www.pygame.org/docs/},
diakses 2026.

\bibitem{sweigart}
Al Sweigart. \emph{Making Games with Python and Pygame}. 2nd edition,
2012. Rujukan teknik, bukan sumber teks.

\end{thebibliography}
